% Generated draft: requires mathematical and source review.
\documentclass[11pt]{article}
\usepackage[margin=25mm]{geometry}
\usepackage{fontspec}
\setmainfont{Latin Modern Roman}
\usepackage{amsmath,amssymb,mathtools,unicode-math}
\setmathfont{Latin Modern Math}
\usepackage{enumitem,needspace,xurl,hyperref}
\setlength{\parindent}{0pt}
\setlength{\parskip}{5pt}
\setlength{\emergencystretch}{3em}
\newcommand{\N}{\mathbb{N}}
\newcommand{\Z}{\mathbb{Z}}
\newcommand{\Q}{\mathbb{Q}}
\newcommand{\R}{\mathbb{R}}
\newcommand{\C}{\mathbb{C}}
\newcommand{\F}{\mathbb{F}}
\newcommand{\A}{\mathbb{A}}
\newcommand{\PP}{\mathbb{P}}
\newcommand{\E}{\mathbb{E}}
\newcommand{\eps}{\varepsilon}
\newcommand{\dd}{\,\mathrm{d}}
\begin{document}
\section*{UnsolvedMath ID: 30001483}
\textbf{Source problem code:} catalogue entry
\subsection{Definitions and mathematical statement}
Let \(h(t)\) denote a solution path for the thin liquid-film evolution intended by the source, let \(E\) denote its energy, and let \(m\) denote its mobility. These are placeholders: the supplied record does not specify the evolution equation, the energy functional, or the admissible solution class. The mobility is required to be superlinear, but the meaning and range of this condition are not specified.

\textbf{Source-level question.} For thin liquid-film evolution with superlinear mobility, obtain estimates of the steepness of \(E\) along solution paths \(h(t)\) sufficient to establish optimal rigorous upper bounds on the coarsening rate.

\textbf{Unresolved formulation.} The record does not define energy steepness, the quantity measuring coarsening, or the criterion of optimality. It also supplies no spatial domain, boundary conditions, initial-data assumptions, parameter ranges, or time regime. Consequently, a precise quantified inequality cannot be extracted without additional source information. In particular, no particular thin-film equation, definition of energy slope, or coarsening exponent is asserted here.
\subsection{Short English statement}
The problem asks whether energy-steepness estimates along thin-film solution paths with superlinear mobility can yield optimal upper bounds on coarsening. The supplied formulation leaves the model and the quantities to be estimated unspecified, so it remains a research question requiring further formulation and human mathematical review.
\subsection{Sources}
\begin{itemize}
\item Local registry: \texttt{data/unsolvedmath/problems.json}, permanent UnsolvedMath ID \(30001483\), source problem code \texttt{catalogue entry}, title \emph{Energy-Slope Bounds for Thin-Film Coarsening}.
\item Supplied report citation: \emph{Phase Transitions} (2010), Oberwolfach Reports, pp.\ 1431--1484, DOI \texttt{10.4171/owr/2010/24}. Source URL and supplied literature source labelled \emph{OWR 2010/24}: \url{https://doi.org/10.4171/owr/2010/24}. The record attributes the question to the report's workshop open-problem discussions; no individual proposer is supplied.
\item Literature assessment dated \texttt{2026-08-21}: the supplied curation labels the problem open, describes the dynamics question as largely open in the source, and reports that its search found no result resolving the extracted broad claim. This is supplied dated curation, not an independently verified literature assessment; time-sensitive status and cited claims require independent verification.
\item The registry labels the statement \emph{exact} and reports that its mathematical content was checked against the official report, with notation, context, and question form normalized without changing the source-level claim. This is a supplied verification claim; the report has not been independently consulted for this draft.
\end{itemize}
\end{document}
